Robust CP methods are heavily dependent on the quality and compatibility of the underlying model. Therefore, we use separate neural networks to ensure a fair evaluation of robust CP methods. In the following sections, we define the models we used for benchmarking both robust CP and worst-case vanilla CP coverage bounds.

\subsection{System Requirements}

All experiments were conducted on a system equipped with two NVIDIA GeForce RTX 4090 GPUs, each providing 24 GB of GDDR6X memory.

\subsection{VRCP Models}

Given that formal verification methods have an inherent trade-off between computational overhead and tightness of their estimations, we use specifically tailored models from Appendix B of \citet{vrcp}. We then compare these models to Lipschitz by design models in the experiments of \S~\ref{sec:resultsclassic} and \S~\ref{sec:resultscoverage}.

Importantly, using larger and deeper models with verification methods is not only computationally punitive, it also usually tends to worsen the obtained bounds as they become looser and looser.

\subsection{Smoothing Models}

We use a ResNet50 architecture in the experiments we run using \emph{CAS} as it offers a good baseline for robust CP results as done in \cite{cas}. Given the already costly nature of smoothing methods at inference, we follow a standard scheme at training time to avoid additional computational burden. An exception stands for the \texttt{CIFAR-10} dataset: models trained with Gaussian noise as per the procedure of \citet{salman2020provablyrobustdeeplearning} still fail in low MC iteration settings. Indeed, we test a ResNet110 model trained with noise augmentation from the aforementioned procedure: yet it still exhibits uninformative set sizes of $10$ for $n_\mathrm{mc}=64$ on every point of the test set. Moreover, we use the recommended $\sigma$ value of $\sigma = 2\cdot \epsilon$ for every smoothing experiment conducted.

\subsection{Lipschitz Models}
\label{app:model_implementation}

The architecture we use for our experiments on \textsl{CIFAR-10}, \textsl{CIFAR-100} and \textsl{TinyImageNet} datasets follows a ResNeXt-like design, with an initial convolutional stem followed by four stages of increasing feature channels (64 → 128 → 256 → 512). Each stage comprises grouped-convolution residual blocks, with downsampling applied between stages to reduce spatial dimensions. The network then employs global average pooling, a fully connected bottleneck, and a classification layer. 

%\hfill

Regarding our experiments on the \textsl{ImageNet} dataset, we leverage a more efficient implementation of Lipschitz parametrized networks, as described in \cite{boissin2025adaptiveorthogonalconvolutionscheme}. This implementation can be found in the \href{https://github.com/deel-ai/orthogonium}{\texttt{orthogonium}} library.

Regarding our \textsl{ImageNet} model we use the following architecture. Starting with two strided convolutions (7×7 → 3×3 kernels) that reduce spatial resolution while expanding channels to 256. Four subsequent processing stages progressively double channel dimensions (256 → 512 → 1024 → 2048), each containing three residual blocks with depthwise 5×5 convolutions. Blocks employ channel expansion/contraction (2 × width), GroupSort2 activations, and learned residual mixing that respects 1-Lipschitz constraints. Then, the spatial features are condensed via global average pooling (7×7 window) before final classification by a linear layer that enforces the 1-Lipschitz condition on each output. 

\paragraph{Training objectives} Our neural networks are trained with Optimal Transport inspired objectives such as the Hinge Kantorovich-Rubinstein loss function introduced in~\cite{acheivingrobustness}. These loss functions are particularly effective to enforce model robustness.

\subsection{Training Hyperparameters}

We train our Lipschitz neural networks with the AdamW optimizer with a learning rate 1e-3. Also, we use the \texttt{SoftHKRMulticlassLoss} from the \href{https://github.com/deel-ai/deel-torchlip}{\texttt{deel-torchlip}} library with the following standard values:

\begin{table}[h]
  \centering
  \sisetup{
    detect-weight = true,
    detect-family = true,
    table-number-alignment = center
  }
  \begin{tabular}{l
                  S[table-format=1.1]
                  S[table-format=1.1]
                  S[table-format=3.0]
                  S[table-format=1.3]}
    \toprule
    \textbf{Dataset}   & {\textbf{Margin}} & {\textbf{Temperature}} & {\textbf{Epochs}} & {\textbf{Alpha}} \\
    \midrule
    \textsl{CIFAR-10}           & 0.6               & 5.0                    & 130               & 0.975            \\
    \textsl{CIFAR-100}          & 0.6               & 5.0                    & 220               & 0.975            \\
    \textsl{TinyImageNet}       & 0.3               & 5.0                    &  80               & 0.975            \\
    \bottomrule
  \end{tabular}
  \caption{Standard training hyperparameter values for the loss function on the different datasets.}
  \label{tab:hyperparams}
\end{table}
